\chapter{Literaturverzeichnis}

Bernd Rechel, E. R. (2014). Trends in health systems in the former
Soviet countries. \emph{European Journal of Public Health}.

Ketevan Glonti, B. R. (2013). Health Targets in the Former Soviet
Countries: Responding to the NCD Challenge? \emph{Public Health Reviews,
Vol. 35, No 1}.

The World Bank. (2016). \emph{SYSTEMATIC COUNTRY DIAGNOSTIC for
UZBEKISTAN.} Washington D.C: The World Bank.

World Health Organization. (2022). \emph{Health Systems in Action
Uzbekistan.} World Health Organization.

Mayring, P. (2015). \emph{Qualitative Inhaltsanalyse: Grundlagen und Techniken}
(12., überarbeitete Auflage). Beltz.

Bashshur, R. L., Shannon, G. W., Smith, B. R., Alverson, D. C.,
Antoniotti, N., Barsan, W. G., ... Krupinski, E. A. (2014). The
empirical foundations of telemedicine interventions in primary care.
\emph{Telemedicine and e-Health, 20}(5), 342--375.

Shigekawa, E., Fix, M., Corbett, G., Roby, D. H., \& Coffman, J.
(2018). The current state of telehealth evidence: A rapid review.
\emph{Health Affairs, 37}(12), 1975--1982.

Ekeland, A. G., Bowes, A., \& Flottorp, S. (2010). Effectiveness of
telemedicine: A systematic review of reviews. \emph{International Journal
of Medical Informatics, 79}(11), 736--771.

Venkatesh, V., Morris, M. G., Davis, G. B., \& Davis, F. D. (2003).
User acceptance of information technology: Toward a unified view.
\emph{MIS Quarterly, 27}(3), 425--478.

Greenhalgh, T., Wherton, J., Papoutsi, C., Lynch, J., Hughes, G.,
A'Court, C., ... Shaw, S. (2017). Beyond adoption: A new framework for
theorizing and evaluating nonadoption, abandonment, scale-up, spread, and
sustainability of health and care technologies (NASSS). \emph{Journal of
Medical Internet Research, 19}(11), e367.

Brooke, J. (1996). SUS: A ``quick and dirty'' usability scale. In
P. W. Jordan, B. Thomas, B. A. Weerdmeester, \& I. L. McClelland
(Eds.), \emph{Usability Evaluation in Industry} (pp. 189--194).
Taylor \& Francis.

Nielsen, J. (1994). \emph{Usability Engineering}. Academic Press.

ISO. (2018). \emph{ISO 9241-11:2018 Ergonomics of human-system
interaction---Part 11: Usability: Definitions and concepts}. International
Organization for Standardization.

Maramba, I., Chatterjee, A., \& Newman, C. (2019). Methods of usability
testing in the development of eHealth applications: A scoping review.
\emph{International Journal of Medical Informatics, 126}, 95--104.

Kruse, C. S., Frederick, B., Jacobson, T., \& Monticone, D. K. (2017).
Cybersecurity in healthcare: A systematic review of modern threats and
trends. \emph{Technology and Health Care, 25}(1), 1--10.

Appari, A., \& Johnson, M. E. (2010). Information security and privacy in
healthcare: Current state of research. \emph{International Journal of
Internet and Enterprise Management, 6}(4), 279--314.

Bass, L., Clements, P., \& Kazman, R. (2021). \emph{Software Architecture
in Practice} (4th ed.). Addison-Wesley.

Richards, M., \& Ford, N. (2020). \emph{Fundamentals of Software
Architecture: An Engineering Approach}. O'Reilly Media.

Humble, J., \& Farley, D. (2010). \emph{Continuous Delivery: Reliable
Software Releases through Build, Test, and Deployment Automation}.
Addison-Wesley.

Beyer, B., Jones, C., Petoff, J., \& Murphy, N. R. (Eds.). (2016).
\emph{Site Reliability Engineering: How Google Runs Production Systems}.
O'Reilly Media.

\chapter{Anhang}

Der Anhang enthält ergänzende Materialien, die im Haupttext referenziert
werden, darunter zusätzliche Dokumentations- und Auswertungsunterlagen.
Er dient der Nachvollziehbarkeit, ohne den Fließtext mit ausführlichen
Detailnachweisen zu überladen.

\chapter*{Eidesstattliche Erklärung}
\addcontentsline{toc}{chapter}{Eidesstattliche Erklärung}

Hiermit erkläre ich, dass ich die vorliegende Arbeit selbstständig
verfasst habe, dass ich sie zuvor an keiner anderen Hochschule und in
keinem anderen Studiengang als Prüfungsleistung eingereicht habe und
dass ich keine anderen als die angegebenen Quellen und Hilfsmittel
benutzt habe. Alle Stellen der Arbeit, die wörtlich oder sinngemäß aus
Veröffentlichungen oder aus anderweitigen fremden Äußerungen entnommen
wurden, sind als solche kenntlich gemacht.

Diese Erklärung erfolgt gemäß § 60 Abs. 8 AllgStuPO der Technischen
Universität Berlin.

Ort, Datum Unterschrift

